\section{第一层因果识别的完整诊断}\label{app:identification}

\textbf{本附录的目的不是报告一个估计，而是报告为什么报告不出一个估计。} 正文第~\ref{sec:identification}~节给出诊断与结论；本附录给出全部数字与复现命令。凡本附录中的数字，均由当前仓库（\texttt{panel\_cityyear.csv} 与 \texttt{policy.csv}）重新跑出；旧面板时代的数字已单独标注并说明为何不再适用。

\subsection{样本、批次与处理变量覆盖}

三个政策批次，方向不一：

\begin{table}[H]
\centering
\caption{三个政策批次}
\label{tab:batches-a}
\small
\begin{tabular}{L{1.2cm}L{6.6cm}L{2.2cm}L{2.4cm}R{1.5cm}}
\toprule
\textbf{批次} & \textbf{城市} & \textbf{方向} & \textbf{同省对照} & \textbf{处理单位} \\
\midrule
2023 & 信阳、南阳、平顶山、洛阳、郑州 & 退出 & 河南其余 12 城 & 5 \\
2024 & 东营、滨州 & \textbf{进入} & 山东其余城市 & 2 \\
2026 & 临汾、吕梁、大同、太原、晋中、晋城、朔州、运城、长治、阳泉 & 退出（省级） & \textbf{0} & \textbf{1} \\
\bottomrule
\end{tabular}
\end{table}

\textbf{处理变量的覆盖：} 3,672 个城市-年中 \texttt{policy\_type} 非空 243 个（6.6\%）；306 城中仅 34 城有任何政策记录，其余 272 城全空；估计窗口 2021—2026 的 1,836 行中非空 201 行（10.9\%）。空白按 0 处理是记录口径，不是事实。

\begin{quote}
\textbf{一处更正（2026-10-20）。} 本附录上一版写“样本期内没有任何一次限放$\rightarrow$禁放的切换。唯一一次是武汉 2006$\rightarrow$2013”。\textbf{这是错的。} \texttt{policy.csv} 中东营（2024-01-02 施行，政府令第 208 号，\texttt{verified=Y}）与滨州（2024-01-01，\texttt{verified=N}）都是限$\rightarrow$禁；面板已将两城 2024—2026 的 \texttt{ban} 置 1。上一版的批次表只有 2023 河南与 2026 山西，漏掉了 2024 山东。现已补入。
\end{quote}

\subsection{诊断零：安慰剂结果变量}

用节前基线作结果变量。燃放管制不应改变年前的浓度。

\begin{table}[H]
\centering
\caption{安慰剂结果变量}
\label{tab:placebo-outcome-a}
\small
\begin{tabular}{L{4.4cm}R{2.0cm}R{2.2cm}R{2.4cm}}
\toprule
\textbf{设定} & $\beta$ & \textbf{SE} & \textbf{有效聚类数} \\
\midrule
城市 $+$ 年份 & $+10.62$ & 2.54 & 306 \\
城市 $+$ 省份 $\times$ 年份 & $-4.48$ & 1.59 & 33 \\
\bottomrule
\end{tabular}
\end{table}

两个设定都显著且方向相反。基线被处理变量预测，平行趋势不成立。

\subsection{诊断一：安慰剂“日”扫描}

\begin{table}[H]
\centering
\caption{安慰剂“日”扫描}
\label{tab:placebo-days-a}
\small
\begin{tabular}{lR{1.3cm}R{1.3cm}R{1.6cm}R{1.3cm}R{1.3cm}}
\toprule
\multicolumn{6}{l}{\textbf{河南（同省对照，17 城）}} \\
\texttt{rel\_day} & $-6$ & $-4$ & $\mathbf{+1}$ & $+3$ & $+6$ \\
$\beta$ & $+18.7$ & $-15.0$ & $\mathbf{-10.8}$ & $+35.3$ & $-14.2$ \\
$t$ & $+2.38$ & $-2.60$ & $\mathbf{-0.56}$ & $+3.21$ & $-3.15$ \\
SE & 7.9 & 5.8 & 19.3 & 11.0 & 4.5 \\
\midrule
\multicolumn{6}{l}{\textbf{山西 2026（全国对照）}} \\
\texttt{rel\_day} & $-7$ & $-4$ & $\mathbf{+1}$ & $+5$ & $+6$ \\
$\beta$ & $-8.4$ & $-28.0$ & $\mathbf{+37.7}$ & $+24.9$ & $+54.6$ \\
$t$ & $-2.80$ & $-8.03$ & $\mathbf{+3.02}$ & $+2.81$ & $+4.49$ \\
\bottomrule
\end{tabular}
\end{table}

河南的初一不显著（$t=-0.56$），无关日 $-6$、$-4$、$+3$、$+6$ 显著。山西的初一显著（$+37.7$，$t=3.02$），但 $-7$ 到 $-1$ 的每一天也都显著（$t$ 在 $-2.8$ 到 $-8.0$ 之间）——处理组在政策前已在相对下行。

复现：
{\small
\begin{verbatim}
python3 -I scripts/placebo_days.py --province 河南 --treat-year 2023
python3 -I scripts/placebo_days.py --province 山西 --treat-year 2026 \
    --control all
\end{verbatim}
}

\subsection{诊断二：事件路径的前置期}

\begin{table}[H]
\centering
\caption{各批次事件路径（原始差）}
\label{tab:preperiods-a}
\small
\begin{tabular}{L{2.0cm}L{8.4cm}R{2.6cm}}
\toprule
\textbf{批次} & \textbf{前置期} & \textbf{切换当期} \\
\midrule
2023 河南 & $-1$：$-46.61$ & $+75.30$ \\
2024 山东 & $-2$：$-56.87$；$-1$：$+141.97$ & $-81.48$ \\
2026 山西 & $-4$：$-45.83$；$-3$：$+32.42$；$-2$：$+15.24$；$-1$：$-22.33$ & $+44.61$ \\
\bottomrule
\end{tabular}
\end{table}

\noindent 路径是原始差（处理组变化减对照组变化），未按方向翻符号。三个批次的前置期都有与效应同量级的波动。

\subsection{固定效应设定与窗口敏感性}

\begin{table}[H]
\centering
\caption{固定效应设定}
\label{tab:fe-a}
\small
\begin{tabular}{L{4.4cm}R{2.0cm}R{2.2cm}R{2.2cm}R{2.4cm}}
\toprule
\textbf{设定} & $\beta$ & \textbf{SE（CR1）} & \textbf{SE（含校正）} & \textbf{有效聚类数} \\
\midrule
城市 $+$ 年份 & $-31.74$ & 10.60 & 11.6—12.4 & 306 \\
城市 $+$ 省份 $\times$ 年份 & $-5.34$ & 14.41 & 14.8—15.2 & \textbf{33} \\
\bottomrule
\end{tabular}
\end{table}

\begin{table}[H]
\centering
\caption{窗口敏感性（\texttt{--min-year}）}
\label{tab:window}
\small
\begin{tabular}{R{1.9cm}R{1.2cm}R{2.2cm}R{2.9cm}L{4.5cm}}
\toprule
\texttt{min\_year} & \textbf{年数} & \textbf{城市$+$年份} & \textbf{城市$+$省份$\times$年份} & \textbf{合法？} \\
\midrule
2019 & 8 & $-40.62$ & $-8.09$ & $\times$ 河南对照城出现 2020 假切换 \\
2021 & 6 & $-31.74$ & $-5.34$ & $\checkmark$ \\
2022 & 5 & $-40.48$ & $-7.62$ & $\checkmark$ \\
\bottomrule
\end{tabular}
\end{table}

\begin{quote}
\textbf{一处撤回（2026-10-20）。} 上一版报告 $-42.60$（9.52）、$+10.80$（18.77，有效聚类 17），并据此认为“符号随固定效应设定翻转”。那是政策库纳入 2024 山东批次之前的面板。当前面板下两个设定都是负号，\textbf{翻转不成立，该诊断撤回}。复核过的六个数字随之失效，见第~\ref{app:replication}~小节。
\end{quote}

\subsection{精确随机化推断}

河南 17 城取 5，$C(17,5)=6{,}188$，全部穷举。

\begin{table}[H]
\centering
\caption{精确随机化推断}
\label{tab:ri}
\small
\begin{tabular}{L{4.6cm}R{2.2cm}R{2.0cm}R{2.2cm}R{2.2cm}}
\toprule
\textbf{口径} & \textbf{观测统计量} & \textbf{双侧 $p$} & \textbf{位于零分布} & \textbf{零分布 sd} \\
\midrule
初一日均 \pmm{2.5} & $-26.20$ & 0.1836 & $-1.35$ SD & 19.39 \\
\bottomrule
\end{tabular}
\end{table}

\paragraph{校准性。} 对河南 8 次无政策变动的转移逐次检验，得 $p=0.043$、$0.173$、$0.447$、$0.516$、$0.587$、$0.850$、$0.929$、$0.999$，KS 距离 0.225（$n=8$ 的 5\% 临界值约 0.457，不拒绝）。

\paragraph{边界。} 本检验不是因果零假设下的检验；政策内生正是主要担忧。它真正有价值的产出是噪声尺度。

\subsection{逐城贡献与方向异常}

\begin{table}[H]
\centering
\caption{切换当期逐城与对照之差（\ugm，原始差）}
\label{tab:per-city}
\small
\begin{tabular}{L{2.4cm}L{10.6cm}}
\toprule
\textbf{批次} & \textbf{逐城} \\
\midrule
2023 河南（退出） & 信阳 $+36.4$；平顶山 $+47.4$；洛阳 $+74.2$；南阳 $+88.6$；郑州 $+129.8$ \\
2024 山东（进入） & 滨州 $-120.7$；东营 $-42.3$ \\
2026 山西（退出） & 太原 $-103.8$；大同 $+0.1$；阳泉 $+31.3$；吕梁 $+36.4$；长治 $+50.3$；运城 $+62.7$；晋城 $+74.3$；晋中 $+90.7$；临汾 $+96.1$；朔州 $+107.6$ \\
\bottomrule
\end{tabular}
\end{table}

2023 批次五城全部为正（放开后相对上升），方向与理论一致；2024 批次两城都为负（收紧后相对下降），方向也一致。\textbf{方向一致而前置期不平稳，正是“‘看起来对’但不可信”的典型形态。}

山西的太原仍为唯一反号城市且幅度最大（$-103.8$），单城把 10 城均值从 $+60.2$ 拉到 $+43.7$。

\subsection{结构性限制：省级切换是最难识别的变异}

山西全部面板城市在 2026 年退出处理。三条彼此独立的原因令其不可识别：

\begin{enumerate}
  \item \textbf{省内没有对照。} 同省对照库为空——省级通告一夜之间改变全省所有地级市。
  \item \textbf{省份 $\times$ 年份固定效应将其整体吸收。} 处理变异与省份 $\times$ 年份虚拟变量完全共线，对 $\beta$ 的贡献为零。
  \item \textbf{有效处理聚类数是 1，不是 10。} 这些城市来自同一个省政府的同一个决定。城市层聚类给 SE $=14.8$（$t=3.54$），省级聚类给 SE $=1.38$（$t=37.9$，荒谬）。
\end{enumerate}

\textbf{补充证据：山西 vs.\ 邻省（河北、河南、陕西、内蒙古）的逐年差。}

\begin{table}[H]
\centering
\caption{山西减邻省，逐年（当前面板）}
\label{tab:shanxi-neighbor}
\scriptsize
\begin{tabular}{lrrrrrrrrrrrr}
\toprule
年 & 2015 & 2016 & \textbf{2017} & 2018 & 2019 & \textbf{2020} & 2021 & 2022 & 2023 & 2024 & 2025 & \textbf{2026} \\
\midrule
差 & $-4.3$ & $-7.4$ & $\mathbf{+26.0}$ & $-1.7$ & $-33.3$ & $\mathbf{+56.8}$ & $-31.6$ & $-12.4$ & $-34.2$ & $-10.0$ & $-37.1$ & $\mathbf{+53.5}$ \\
\bottomrule
\end{tabular}
\end{table}

2026 年放开的差为 $+53.5$，而 \textbf{2020 年（无任何政策变化）为 $+56.8$，比它还大}。那个“效应”落在历史波动范围之内。

\begin{quote}
\textbf{这一条推翻了本组早先的判断。} 原设想写“省级切换的好处：一夜之间改变全省十几个地级市，是最接近外生冲击的变异”。事实相反：它同时清空了省内对照、被省份 $\times$ 年份虚拟变量吸收、且只有 1 个有效处理聚类。\textbf{它是本文最难识别的变异类型。}
\end{quote}

\subsection{噪声尺度与功效}

\begin{table}[H]
\centering
\caption{噪声尺度}
\label{tab:sigma}
\small
\begin{tabular}{L{9.0cm}R{3.6cm}}
\toprule
\textbf{量} & \textbf{数值} \\
\midrule
$\Delta$(初一 \pmm{2.5}) 城际标准差（306 城 $\times$ 11 次年度转移 $=3{,}366$ 对） & \textbf{84.40} \\
$\Delta$(基线 \pmm{2.5}) 城际标准差 & 32.3 \\
$\Delta$(初一) 与 $\Delta$(基线) 的相关 & $+0.098$ \\
$\Delta$(初一) 的省内标准差（河南，2022$\rightarrow$2023） & \textbf{35.3} \\
同上，跨 9 次转移的范围 & \textbf{35—107} \\
尖峰增量：均值 / 中位数（平衡面板） & $+43.1$ / \textbf{28.1} \\
随机化零分布标准差 & 19.39 \\
\bottomrule
\end{tabular}
\end{table}

\begin{quote}
\textbf{更正：$\sigma$ 是 35.3，不是 46。} 早先的 $\sigma\approx46$ 漏掉了有限总体校正中的 $-2\operatorname{Cov}$ 项。正确公式为
\[
\operatorname{Var}(\text{stat})=\frac{S^2N^2}{(N-1)\,k\,(N-k)},\qquad N=17,\ k=5\ \Rightarrow\ \mathrm{sd}(\text{stat})=0.5487\,S.
\]
但\textbf{不能据此把采集目标降到 $K\approx10$}：$\sigma$ 在转移之间波动 3 倍（35—107），35.3 恰是最有利的那一年。\textbf{规划要用中心值 $\sigma\approx70$。}
\end{quote}

\begin{table}[H]
\centering
\caption{采集目标与功效}
\label{tab:power}
\small
\begin{tabular}{R{1.3cm}R{1.3cm}R{2.8cm}R{2.6cm}L{3.6cm}}
\toprule
$K$ & $M$ & $\sigma=35.3$ & $\sigma=70$ & \textbf{能否检测 43} \\
\midrule
\textbf{5（现状）} & 12 & 53 & 100 & $\times$ \\
10 & 30 & 36 & 68 & 边缘 / $\times$ \\
\textbf{20} & 60 & 26 & \textbf{49} & $\checkmark$（勉强） \\
\textbf{30} & 90 & 21 & \textbf{40} & $\checkmark$ \\
\bottomrule
\end{tabular}
\end{table}

目标为 20—30 个“带同省对照的市级切换城市”，现有 7 个（河南 5、山东 2）。\textbf{但这 20—30 个城市可能根本不存在}——2017—2019 年那波禁放很多是省级统一部署，结构上不可用。

\subsection{三轮独立复核}\label{app:replication}

2026-10-15 安排了三个彼此独立的复核者。\textbf{独立性检验的做法}：复算者与独立分析者被禁止读取结论文档；复算者另被禁止读取或运行 \texttt{scripts/} 下任何文件，必须自写代码。

\begin{table}[H]
\centering
\caption{三轮独立复核（旧面板）}
\label{tab:replication-a}
\small
\begin{tabular}{L{2.0cm}L{3.8cm}L{9.0cm}}
\toprule
\textbf{复核者} & \textbf{独立性} & \textbf{结果} \\
\midrule
A 独立复算 & 未读脚本与结论文档，自写代码，并用全虚拟变量 OLS 第二条路径 & 在当时面板上六个数字逐位一致 \\
B 敌对代码审查 & 读代码，目标是推翻结论 & 七条结论均未被证伪；找到 4 处真 bug 与 2 处文档错；验证交替投影两轮即精确收敛、\texttt{rel\_day} 映射正确、政策库展开正确 \\
C 另起炉灶 & 未读脚本与结论文档，自行重建时间轴与估计量 & 独立得出同一结论；两个结果变量与面板 3672/3672 全部一致 \\
\bottomrule
\end{tabular}
\end{table}

\begin{quote}
\textbf{本小节的适用范围必须限定。} 上述复核在 \textbf{2026-10-15 的面板}上完成，当时政策库尚未纳入 2024 年山东批次。纳入后主估计由 $-42.60$ 变为 $-31.74$，\textbf{复核过的具体数值不再适用于本版}。复核仍然有效的部分，是它对管线的验证：交替投影、\texttt{rel\_day} 映射、政策库展开三项均无 bug，因此“零结果是管线造成的”这一替代解释仍被排除。本版所有估计均由当前面板重新跑出。
\end{quote}

\subsection{已知口径问题}

\begin{table}[H]
\centering
\caption{代码与口径的已知问题}
\label{tab:code-issues}
\small
\begin{tabular}{R{0.7cm}L{3.2cm}L{8.0cm}L{2.4cm}}
\toprule
\textbf{编号} & \textbf{位置} & \textbf{问题} & \textbf{状态} \\
\midrule
1 & 对照均值 & 恒用全体未变动城市，无视同省参数 & 已修 \\
2 & 聚类标准误 & 只乘 $G/(G-1)$，未乘 $(N-1)/(N-K)$ & 已披露；两种口径均报 \\
3 & 聚类数 & 报名义聚类数，未报有效聚类数 & 已修 \\
4 & 窗口截断 & 无法用命令行关闭 & 已修 \\
5 & 随机化脚本 & 只校验首个城市的年份完整性 & 已修 \\
6 & 随机化脚本 & 接受处理定义参数但忽略它 & \textbf{未修} \\
7 & 输出语句 & 一处 \texttt{\%} 未转义，运行到标准误提示时抛 \texttt{TypeError} & \textbf{已修（2026-10-20）} \\
\bottomrule
\end{tabular}
\end{table}

\textbf{第 7 条值得单独说明}：该 bug 使 \texttt{run\_did.py} 在打印完结果后崩溃，因此“数字跑出来了”与“脚本正常结束”是两件事。本版已修，并确认全流程正常退出。

\subsection{本附录的小结：哪些数字可以引用}

\textbf{可以引用：} 三个批次的构成；处理变量覆盖 6.6\%（窗口内 10.9\%）；安慰剂结果变量的两个系数；安慰剂“日”扫描的方向性结论；三个批次的前置期；两个固定效应设定下的 $\beta$ 与有效聚类数（306 与 33）；随机化推断的 $p$ 值与边界；$\sigma=35.3$。

\textbf{不可以引用：}
\begin{itemize}
  \item 任何单一的因果效应点估计。\textbf{本文不报告禁放效应的大小。}
  \item 旧面板的 $-42.60$、$+10.80$ 与“符号翻转”。它们已被当前面板推翻。
  \item “$p=0.85$ 说明检验校准”这一论证。它已被 KS 检验取代。
\end{itemize}
